\section{来源审计：把敏感议题改写成可验证的系统课}

本讲由 Thorn CEO Julie Cordua 主讲，讨论在线儿童安全技术如何帮助平台识别风险、减少重复传播、组织人工复核，并把符合条件的线索交给法定报告与调查体系。公开视频在 2026 年 8 月 11 日访问时要求登录；仓库保留官方封面与 833 条时间戳字幕，因此本文以课堂字幕为叙事主线，以 Thorn、NCMEC 与 NIST 的官方资料校验持久机制，不把后续产品能力倒写成 2025 年课堂事实。

\begin{warningbox}{敏感内容处理与证据边界}
本文不复述、描绘或展示违法内容，也不提供生成、传播或规避检测的方法。课堂中的个案仅抽象为系统事件，例如“平台检测到高风险文件并完成报告”。报告数、文件数、研究比例和产品规模都依赖年份、定义与去重方式；除非明确标注，不能把不同年度数字直接相减或比较。
\end{warningbox}

\begin{importantbox}{本讲的三个闭环}
\begin{enumerate}
  \item \textbf{Response loop}：上传或消息 $\rightarrow$ 检测 $\rightarrow$ 优先级 $\rightarrow$ 人工复核 $\rightarrow$ 平台处置与法定报告；
  \item \textbf{Learning loop}：经授权的已验证结果 $\rightarrow$ trusted data/hash/model update $\rightarrow$ 覆盖未来风险；
  \item \textbf{Governance loop}：新技术与新滥用 $\rightarrow$ threat research $\rightarrow$ Safety by Design $\rightarrow$ 测试、审计与透明度。
\end{enumerate}
\end{importantbox}

\begin{knowledgebox}{术语消化：不要把四种对象混成“坏内容”}
\textbf{Content} 是平台处理的文件或消息；\textbf{signal} 是 hash match、模型 score 或行为模式；\textbf{case} 是按账户、时间和上下文组织的可复核事件；\textbf{outcome} 是平台处置、法定报告、调查或受害者保护结果。工程指标必须说明自己测量哪一层，否则“检测量上升”既可能代表覆盖改善，也可能只是重复文件、阈值变化或报告口径变化。
\end{knowledgebox}

\subsection{为什么这不是“训练一个分类器”}

本节先确定工程对象。\term{CSAM} 是 child sexual abuse material 的缩写，指法律和平台政策所界定的儿童性虐待材料；术语本身描述受害与证据，不应被误写成普通“成人内容”。一个模型只能产生 match 或 score，真正的保护结果还依赖平台 policy、审核队列、证据保全、法定报告机构、执法调查与受害者识别。任何一环失效，前一环的准确率都不能自动转化为安全结果。

\lecturefigure{01-response-pipeline.png}{儿童安全技术是 detect、review、action/report 与 victim identification 组成的端到端响应链。}{本地字幕 00:00--05:50；基于课堂流程的非敏感概念重绘。}

\begin{knowledgebox}{读图：先区分 signal、decision 与 outcome}
左侧 classifier 或 hash service 产生的是\textbf{信号}；平台审核与 policy 形成\textbf{产品决策}；报告机构和调查人员在法定程序中追求\textbf{调查结果}。模型 score 不能直接证明身份、主观意图或法律责任。系统设计必须记录每次状态转换由谁执行、依据什么 policy、留下什么 audit evidence，以及错误如何申诉或纠正。
\end{knowledgebox}

\teachervoice{Cordua 的开场并不是用冲击性案例吸引注意，而是追问：当平台报告负担持续增长时，技术能否让检测、处置与调查人员的有限时间真正连接起来。这个问题把课堂从“坏内容识别”提升为 end-to-end operations。}

\subsection{本章小结}

本讲的目标是解释完整响应系统，而不是给某个模型做产品宣传。后文所有机制都要回答三个问题：它减少了哪类风险、把工作交给了谁、其证据边界在哪里。

